\section{Conclusion}
\label{sec:conclusion}

We have presented \methodname, a topology-aware process reward framework that treats sequential reasoning traces as implicitly structured graphs.
By extracting dependency DAGs from chain-of-thought outputs and computing two deterministic process signals---topological structure reward for global graph validity and continuity reward for local step traceability---\methodname provides dense, verifiable process supervision without learned reward models or human annotations.
The hierarchical reward aggregation (Eq.~\ref{eq:r_total}) preserves answer-correctness primacy while amplifying structural quality, and SCAE prevents process rewards from favouring well-structured but incorrect traces.
On private mathematical critique benchmarks, TopoPRM improves overall accuracy by $+13.4$ points over outcome-only GRPO while reducing average response length by 11\%.
The topology-verified self-distillation stage (TVSD) further compresses the teacher into compact students ($\leq$4B), retaining structural reasoning profiles while halving the generation budget.

\paragraph{Limitations.}
Our DAG extraction relies on rule-based parsing of mathematical expressions and textual markers, which may miss implicit logical dependencies or fail on non-standard notation.
The topology and continuity rewards capture structural properties but do not verify the mathematical correctness of individual derivation steps; semantic verification remains the responsibility of the outcome reward.
The PRM quality validation (Section~\ref{sec:rq2}) reveals that the current topology reward has limited sensitivity to step-order permutations (good vs.\ shuffled win rate $\approx$0.52), suggesting room for improvement in dependency-edge construction.
Finally, the DAG extraction rules are designed for mathematical reasoning; extending to open-domain reasoning would require domain-adapted parsers.

\paragraph{Future Work.}
Several directions merit further investigation:
(i)~combining deterministic topology rewards with generative PRMs~\citep{genprm2026,thinkprm2025} to capture both structural and semantic step quality;
(ii)~extending DAG extraction to multi-modal reasoning traces that combine text, code, and diagrams;
(iii)~exploring on-policy distillation~\citep{agarwal2024gkd,opsdc2026} as an alternative to reverse-KL for the compression stage, using DAG diagnostics as the densification signal;
(iv)~applying the topological modeling and layer-compression pipeline to reasoning tasks beyond mathematics, such as scientific problem solving and code generation.
